\subsection*{12. The Price Waterfall: List Price Is Not Cash Retained}

Imagine a course advertised at ₹50,000. A ₹5,000 promotion makes the pre-tax transaction price ₹45,000. Expected refunds and collections adjustments of ₹2,000 leave ₹43,000 of simplified retained fee. Suppose enrollment acquisition costs ₹7,000, incremental teaching and assessment ₹9,000, and payment and support costs ₹1,000. The illustrative transaction contributes ₹26,000 before fixed costs.

\begin{table}[htbp]
\centering
\small
\begin{tabularx}{\columnwidth}{Xr}
\toprule
\textbf{Stage} & \textbf{Amount} \\
\midrule
Public list price & ₹50,000 \\
Promotional concession & $-$₹5,000 \\
Actual pre-tax fee & ₹45,000 \\
Expected refunds / uncollected amounts & $-$₹2,000 \\
Simplified net realized fee & ₹43,000 \\
Acquisition expense & $-$₹7,000 \\
Incremental instruction and assessment & $-$₹9,000 \\
Payment and variable support & $-$₹1,000 \\
\midrule
\textbf{Contribution before fixed costs} & \textbf{₹26,000} \\
\bottomrule
\end{tabularx}
\caption{Illustrative economic price waterfall from list price to realized contribution.}
\label{tab:price-waterfall}
\end{table}

This is an \textbf{economic pricing illustration}, not a formal revenue-recognition schedule. Actual accounting treatment of consideration, refunds, taxes, and bad debts follows applicable standards and contract terms; Chapter 17 addresses that distinction. The waterfall is useful because a founder can mistake an impressive headline fee for money genuinely retained.

\subsection*{13. The Constraint May Be Expert Time, Not Seats}

Suppose 1 client pays ₹1 lakh but consumes 100 skilled hours, while another pays ₹60,000 and uses 20. After appropriate variable expenses, the 2nd can produce more contribution per constrained hour. In education, a highly priced learner who requires extraordinary custom support can be less economically attractive than a well-matched learner who succeeds through a designed pathway.

A responsible price should be compatible with both the customer promise and the organization's available capacity. Scarce expert attention can justify real premiums, but ``unlimited mentorship'' is dangerous when the availability of a qualified mentor is finite. The bottleneck is sometimes grading, sometimes coaching, sometimes implementation, and sometimes customer acquisition: not the number of course recordings.

\subsection*{14. Cash Timing Is Part of the Price Architecture}

₹10 lakh paid upfront is not economically identical to ₹10 lakh collected in 12 installments with delayed payments and default risk. A prepaid annual contract creates cash today and service obligations tomorrow. An enterprise project with excellent contribution may fail if the company must pay contractors months before the invoice clears.

Set deposit, milestone, acceptance, cancellation, renewal, and credit terms as carefully as the headline fee. Be particularly cautious with expensive consumer education and financing arrangements: the goal is not to hide affordability constraints behind monthly payment language. Transparent cash terms protect the buyer and the seller. A business that can sell profitably but cannot finance fulfillment is not yet financially sound.
